% !TeX root = ../main.tex

\section{Conclusions}
\label{sec:discussion}

This paper proposes a Bayesian counterpart of the ADMM quantile trend filtering estimator: an asymmetric Laplace working likelihood and a Laplace prior with a global rate on the differences of the trend, with the penalty and the likelihood scale learned from the data. The posterior is sampled by a Gibbs sampler whose trend update is a P-MALA step on the whole conditional of the trend; its proximal map is a weighted trend filtering problem, every linear system is banded, and the cost per sweep grows linearly with $n$. For several quantile levels, the joint sampler restricts each trend to the box defined by its neighbors, and every draw is ordered. Theorem~\ref{thm:main} shows that, for a fixed scale and a fixed penalty of order $n^{(2k+1)/2}$, the posterior contracts around the true quantiles at the rate $n^{-(2k+1)/(2k+2)}$, and that this rate is exact. The result holds under two conditions on the distribution functions of the data near the quantile and does not assume the working likelihood.

The numerical results support the method. On the two benchmark signals, Bayesian-PMALA is competitive with the horseshoe sampler of \citet{onizuka2024fast} and with the ADMM estimator of \citet{brantley2020baseline}, and at extreme upper quantiles the fits of the joint sampler never cross. In the peak-detection design it performs close to ADMM, and its time per sweep grows linearly with $n$. In the two data applications, the same sampler estimates the baseline of an air-quality series with $n = 8{,}640$ observations and the growth centiles of Dutch children. A natural next step is a contraction result for the joint posterior of several quantile levels under the ordering constraint, which would extend Theorem~\ref{thm:main} to the non-crossing sampler.

% ----------------------------------------------------------------------------
% PARKED (W-04, 2026-11-24): the Discussion before the rewrite. The Conclusions above replace it.
% Possible homes: computational comparison -> Section 4 (computational cost); global penalty -> Section 4
% (peak detection); unknown sigma -> Section 3.3; joint target -> Section 2 (joint sampler).
% No home yet: the coverage remark and the list of extensions.
% Bayesian-PMALA brings exact posterior sampling to quantile trend filtering while keeping the proximal structure of the optimization problem. Under the normal--exponential representation, the proximal map of the whole conditional posterior of the trend is a weighted trend filtering problem, so no part of the check loss is smoothed and the Metropolis--Hastings correction is exact. In the simulation study each competitor is most accurate on the benchmark that suits it, and Bayesian-PMALA is competitive with both on both signals. It matches the locally adaptive horseshoe sampler with a single global penalty per quantile, and it matches ADMM while learning the penalty and the likelihood scale from the posterior instead of selecting a penalty by an information criterion.
%
% The computational structure is the second difference from the Gibbs sampler of \citet{onizuka2024fast}. Its supplied implementation forms a dense inverse and uses dense Gaussian simulation at each iteration, also with its sparse option. Every linear system in Bayesian-PMALA is banded, with bandwidth $k+1$ for every $n$. The metric $M$ and the ADMM system matrix change with $(z, \lambda, \sigma)$ and are therefore factorized at each sweep, but one factorization of each serves the forward and reverse proximal maps and the proposal noise, and for a fixed ADMM iteration budget the work per sweep grows linearly with $n$. The measured running times differ accordingly (Section~\ref{sec:scalability}): a sweep of the supplied implementation takes about $75$ times as long as a Bayesian-PMALA sweep at $n = 500$ and about $7\,000$ times as long at $n = 4000$. Running time per effective posterior draw also depends on mixing, which we do not compare here.
%
% The experiment of Section~\ref{sec:air_simulation} shows what the global penalty buys. There the target is a peak-free baseline beneath transient peaks, so local adaptivity is a liability: the horseshoe's local scales let the fit follow the peaks, whereas a single penalty cannot bend locally. Bayesian-PMALA is far more accurate than the horseshoe sampler on that design and close to ADMM.
%
% Treating the ALD scale $\sigma$ as unknown is a substantive modeling choice. At fixed $(\lambda, \sigma)$ the conditional mode is the quantile trend filter with penalty $\lambda\sigma$ \eqref{eq:map_qtf}, so the posterior learns the effective penalty through two parameters, the scale adapting to the noise level and $\lambda$ to the roughness of the trend, rather than asking one tuning parameter to absorb both roles. Learning $\sigma$ does not by itself guarantee frequentist coverage when the asymmetric Laplace likelihood is misspecified, a point also recognized by the calibration approach of \citet{onizuka2024fast}.
%
% A single global penalty is a deliberate simplification relative to coefficient-specific shrinkage, and the experiments do not establish a universally preferred prior or difference order. The pointwise intervals are below nominal coverage at the outer quantile levels (Section~\ref{sec:sim:ind:smooth}), and their calibration there requires further study.
%
% The joint sampler of Section~\ref{subsec:joint} enforces non-crossing on every posterior draw at the observed indices, and it does so without leaving the P-MALA framework: the order constraint adds the indicator of a box to the negative log-conditional of each trend, which keeps it convex, so only the proximal map changes and the Metropolis--Hastings correction keeps its single-level form. The joint target is a constrained product of working posteriors for the same data, not a joint likelihood for several quantiles, and how the constraint shifts each level's posterior relative to its single-level fit, in location and in interval width, remains to be studied.
%
% Several extensions follow from the present work. Locally adaptive priors can be combined with the same proximal update, since the conditional of the trend remains a weighted quadratic plus a weighted $\ell_1$ penalty. Extension to spatially indexed signals requires generalizing the difference operator to graphs, for which the systems are no longer banded and sparse or multigrid solvers may serve a similar role. Finally, the convergence properties of the full Gibbs sampler with continuous step-size adaptation remain to be established.
% ----------------------------------------------------------------------------
